% ══════════════════════════════════════════════════════════════════════════
% DROP-IN: Insert this section between \section{Introduction} and
%          \section{Data} in draft_informs_2026.tex
%
% Also replace the \todo{2--3 sentences on prior literature...} in the
% Introduction with the two sentences marked [INTRO CITE BRIDGE] below.
%
% INTRO CITE BRIDGE — paste these 2 sentences into the intro TODO slot:
%   Prior research has established that FDA inspection outcomes predict
%   downstream drug shortages \citep{Stomberg2018, Wang2025MSOM} and that
%   machine learning can forecast shortages from pharmacy demand signals
%   \citep{Liu2021, Pall2023}.
%   However, the rich narrative text that inspectors generate in Form 483
%   observations has not been used as a predictive input, and no study has
%   directly linked inspection-text signals to facility-level patient harm
%   measured through adverse event reporting.
% ══════════════════════════════════════════════════════════════════════════

\section{Related Work}
\label{sec:lit}

\subsection{NLP and LLMs for FDA Regulatory Text}

Applying machine learning to FDA regulatory documents has concentrated on two
corpora: drug labels and adverse-event reports.
On drug labels, \citet{Bayer2021} evaluated 23 NLP systems for adverse drug
event extraction and found that the best systems reach F1 $\approx 0.89$,
suitable for human-assisted review but not fully autonomous processing.
\citet{Tatonetti2025OnSIDES} fine-tuned PubMedBERT on 200 curated labels to
produce 7.1 million drug--adverse-event pairs from 47,211 FDA and EMA labels
(F1 $= 0.90$; AUROC $= 0.92$).
\citet{Wu2021BERT} achieved Matthews Correlation Coefficients of 0.84--0.87
on drug-induced liver injury classification from FDA labels, compared to 0.60
for keyword rules.
A common finding across these studies is that fine-tuned domain-adapted models
(PubMedBERT, BioBERT) outperform general-purpose models on tasks requiring
specific biomedical vocabulary---but this advantage has narrowed sharply as
generative LLMs have grown in capability.
\citet{Li2026LLMDevice} demonstrated that GPT-4-class models extract structured
data from FDA medical device adverse-event reports (MDRs) and pre-market
summaries at 93--97\% agreement with human annotators, \emph{without any
fine-tuning on domain corpora}---a result that effectively subsumes the
PubMedBERT fine-tuning advantage for well-specified extraction tasks.
Similarly, \citet{Hassani2025} found GPT-4o achieves 89\% precision and 87\%
recall classifying regulatory provisions in food-safety regulations without
task-specific training.
\citet{Atalay2026RecallRisk} applied multi-task BERT to 54,165 FDA device
recall narratives (0.963 accuracy), confirming that FDA regulatory free text is
amenable to automated classification at scale.

On FDA inspection text specifically, \citet{Pazhayattil2020} conducted the
only published quantitative study of Form 483 data using pre-structured citation
frequency counts from a regulatory database (2014--2018, trend analysis of
observation rates by domain).
Warning-letter text has been analyzed descriptively for domain trends and data
integrity patterns \citep{Rathore2022, Park2025, Kwiecinski2024}, but not used
for feature extraction.
No prior study has applied NLP or LLM methods to the raw narrative text of Form
483 observations to extract structured quality-risk features.

We address this gap.
Consistent with \citet{Li2026LLMDevice}, we apply a modern generative LLM with
structured JSON schema output rather than a fine-tuned BERT model; our semantic
lift analysis (Section~\ref{sec:results}) shows that 27\% and 21\% of
contamination and patient-risk signals, respectively, require contextual
inference that any keyword-based classification approach cannot replicate.

\subsection{Inspection Outcomes and Drug Quality Events}

The empirical relationship between FDA inspection activity and drug supply
disruptions is well-established.
\citet{Stomberg2018} provided the first regression-based evidence that
inspection and citation intensity are significantly associated with drug
shortage rates at the drug level.
\citet{Wang2025MSOM} extended this with instrumental variable analysis on US
manufacturing facilities: OAI inspection outcomes predict a 96.4\% lower
shortage likelihood in the subsequent 12 months, and VAI outcomes predict a
77.2\% reduction.
Interpreted alongside Stomberg, the pattern suggests that inspections trigger
corrective action that ultimately stabilizes supply---but the severity of the
underlying quality problem before inspection is not captured by binary OAI/VAI
classifications.

Machine learning models predict drug shortages from demand-side signals with
substantial accuracy.
\citet{Liu2021} achieved C-statistic $= 0.93$ on pharmacy observations using IV
status, generic-only market structure, and number of manufacturers.
\citet{Pall2023} achieved 69\% accuracy one month ahead using XGBoost on
pharmacy sales records for 784 Canadian drugs.
Both studies relied exclusively on market-structure and demand features; neither
incorporated manufacturing inspection or quality signals.
\citet{Kosmas2023} modeled FDA inspection timing as a POMDP to maximize
expected value of inspecting high-risk facilities; their model assumes quality
is revealed by inspection but cannot observe it from text.

Our paper complements this body of work by providing the supply-side quality
signal these models assume but cannot currently observe: structured features
derived from the narrative content of Form 483 observations.

\subsection{Adverse Events as a Manufacturing Quality Outcome}

FAERS adverse event reports are widely used for pharmacovigilance
disproportionality analysis at the drug level, but their potential as a
facility-level manufacturing quality metric has received limited attention.
\citet{Sardella2021} provided the conceptual foundation: pharmacovigilance is a
primary mechanism for detecting manufacturing-origin safety failures
post-approval, as demonstrated by the nitrosamine contamination crisis and the
heparin adulteration scandal, both of which were first surfaced through adverse
event signal clustering tracing to facility-level manufacturing decisions.
\citet{Brown2020} went further, calling for systematic use of FDA's Sentinel
and FAERS data to generate facility-level manufacturer quality signals, and
noting that the FDA already folds FAERS ``hazard signals'' into its risk-based
facility site-selection model for inspections---the agency itself treats
adverse event patterns as a quality-surveillance signal.

Using FAERS as a manufacturer-level quality comparator introduces potential
confounds.
\citet{Rahman2017FAERS} and \citet{Alatawi2018} demonstrated that FAERS
reporting rates differ between branded and generic products partly due to
perception bias (reporters over-attribute adverse events to generic products),
and developed an authorized-generic comparison design to control for this.
Our analysis largely sidesteps this confound because we compare facilities
manufacturing the \emph{same} API---inter-facility differences within a drug
control for drug-level perception bias.
Consistent with \citet{Potter2025}, we treat FAERS counts as a
\emph{predictive correlate} of manufacturing quality, not as a measure of
adverse event incidence or a causal indicator.

Critically, no prior study has used FEI-aggregated FAERS counts as the
dependent variable in a model of inspection-text quality signals, or validated
FAERS as a quality proxy against independent third-party chemical testing.
We contribute both: a cross-lag analysis of 483-text features against
facility-level FAERS counts at horizons of 0, 1, and 2 years post-inspection,
and external validation of the FAERS quality signal against Valisure
independent analytical testing results for the same 14 APIs.

\subsection{Positioning}

Table~\ref{tab:positioning} summarizes our contribution relative to the closest
prior work.

\begin{table}[h!]
\centering
\small
\caption{Positioning relative to prior work}
\label{tab:positioning}
\begin{tabular}{p{3.5cm} p{2.2cm} p{2.2cm} p{2.2cm} p{2.5cm}}
\toprule
& \textbf{Pazhayattil 2020} & \textbf{Wang et al.\ 2025} & \textbf{Liu/Pall 2021/23} & \textbf{This paper} \\
\midrule
483 observation text  & Counts only   & No  & No  & LLM extraction \\
Inspection outcome    & Citation freq & OAI/VAI  & No  & OAI/VAI + text \\
Quality outcome       & None          & Shortage & Shortage & FAERS + Valisure \\
Facility-level FEI    & Partial       & Yes      & No       & Yes \\
External validation   & No            & IV       & CV       & Valisure testing \\
\bottomrule
\end{tabular}
\end{table}
